% RESEARCH_DRAFT_SHA256: 910b531807b87f8690e0796ba0def13fbebb783f0aae6fbbd6c0e55836289889
% Layout and source-reference wrappers expanded for the web renderer.
\documentclass[11pt]{article}
\usepackage[margin=25mm]{geometry}
\usepackage{fontspec,unicode-math}
\setmainfont{Latin Modern Roman}
\setmathfont{Latin Modern Math}
\usepackage{amsmath,amssymb,mathtools,xurl,hyperref}
\hypersetup{colorlinks=true,urlcolor=blue,linkcolor=blue}
\setcounter{secnumdepth}{0}
\setlength{\parindent}{0pt}
\setlength{\parskip}{5pt}
\setlength{\emergencystretch}{3em}
\newcommand{\sref}[2]{\hyperlink{source:#1:#2}{[S#2]}}
\newcommand{\cataloguescope}[1]{\paragraph{Source record.}#1}
\newcommand{\floor}[1]{\left\lfloor#1\right\rfloor}
\begin{document}
\section*{Research-attempt provenance}
This standalone research attempt retains the complete original definition body below,
including its historical source and preparation notes. Those notes describe the
earlier catalogue preparation; the new attempt and its verification are identified
separately after that body. No original ID, catalogue status, or source statement is changed.

\section{Discrete Vector Fields for Voronoi-Cohomology Computations}\label{statement}

\subsection{Definitions and mathematical statement}

\paragraph{Shared notation.}Standard number systems are $\mathbb N,\mathbb Z,\mathbb Q,\mathbb R,\mathbb C$. Specialized notation and index ranges follow the cited source and the definitions below.


A discrete vector field pairs incident cells, each cell used at most once. An acyclic such matching provides a discrete-Morse reduction that can simplify chain complexes used to compute cohomology.
\paragraph{Record scope and qualifications.}
The catalogue title is retained as metadata. Named certificates, contracts, bounds or frameworks in that title are not asserted as results of the original source extract.
The meaning of effective, admissible symmetry constraints, and computational target are not quantified.

Source-record limitation: The effective application of a technique is a broad methodological prompt without a specified mathematical or runtime guarantee.
\paragraph{Statement recorded in the supplied source.}
Can the discrete vector-field technique be applied effectively to cohomological computations with the Voronoi cell complex? \href{https://doi.org/10.4171/owr/2016/52}{[S2]}
\paragraph{Source status.}
The supplied catalogue labels this source record open; that metadata was checked against the supplied snapshot on 2026-10-11. The label is preserved for provenance and is not an independent verification of present open status or mathematical correctness. The catalogue literature assessment is dated 2026-08-21.
\subsection{Short English statement}
Use discrete vector fields to simplify the cited Voronoi cohomology computations.
\subsection{Sources}
\begin{itemize}
\item[S1] Ulam AI and the UnsolvedMath Contributors, UnsolvedMath: A Curated Collection of Open Mathematics Problems, catalogue entry. Catalogue curation and metadata: CC BY 4.0. The source extract is retained; definitions, English summaries and qualifications are editorial additions.
\url{https://huggingface.co/datasets/ulamai/UnsolvedMath}.
\item[S2] Mini-Workshop: Computations in the Cohomology of Arithmetic Groups (2017), Oberwolfach Reports, DOI 10.4171/owr/2016/52, pp. 2941--2973
\url{https://doi.org/10.4171/owr/2016/52}.
\item[S3] OWR 52 preliminary report
\url{https://orbilu.uni.lu/bitstream/10993/29499/1/preliminary_OWR_2016_52.pdf}.
\item Source-specific Oberwolfach Reports licences, including CC BY-SA 4.0 where applicable, are retained; see THIRD\_PARTY\_NOTICES.md.
\end{itemize}
\paragraph{Preparation and review.}
The original source and UnsolvedMath contributors retain attribution for the source extract. Codex (GPT-6; exact serving revision not exposed) assisted with definitions, English summaries, scope notes and formatting. No new human mathematical authorship or independently verified solution is asserted. Catalogue identity checked against the complete supplied registry and current repository. Distinct upstream records, including repeated questions and broad or damaged extracts, are retained for complete coverage of the source CS category. A repeated record is not a new independent conjecture. Complete mathematical-equivalence and current-literature review has not been performed. Preparation and final review follow contribution prompts C and D, with source gaps disclosed. Structural and rendering checks do not establish a new theorem or open status.
\label{end:statement}

\section{Literature and context check}
Problem 5 in OWR 52/2016, p.~2971, was read: it proposes applying discrete vector fields to Voronoi cohomology computations, without specifying a matching or a guaranteed reduction. The calculation below rederives the algebraic cancellation that such a method must preserve and exhibits two concrete failure modes. It is not an implementation or a new theorem about all arithmetic Voronoi complexes.
\section{Research attempt: Exact chain cancellation and the torsion boundary}

\paragraph{One cancellation, including its correction terms.}
Let $C_*$ be a finite chain complex of finite free modules over a commutative ring. Suppose one differential, in chosen bases, has a unit entry $u$ in the top-left position:
\[
 d_k=\begin{pmatrix}u&a\\b&D\end{pmatrix}.
\]
Write the source basis as $e_1,e_2,\ldots$ and the target basis as $f_1,f_2,\ldots$. Replace
\[
 f_1\text{ by }f'_1=u^{-1}d_ke_1
       =f_1+\sum_{i>1}u^{-1}b_i f_i,
 \qquad
 e_j\text{ by }e'_j=e_j-u^{-1}a_j e_1\quad(j>1).
\]
These are invertible basis changes, and the differential becomes the direct block matrix with blocks $u$ and
\[
 D-bu^{-1}a.
\]
There is no remaining entry in the first row or column outside $u$. The equation $d_{k-1}d_k=0$ gives $d_{k-1}f'_1=0$. The equation $d_kd_{k+1}=0$ gives zero $e_1$ coordinate in $d_{k+1}$, since its first image coordinate is $u$ times that coordinate and $u$ is a unit. Thus the two generators $e_1,f'_1$ split off as a contractible subcomplex, and the remaining differential blocks still square to zero.

The contracting homotopy on this pair is $h(f'_1)=u^{-1}e_1$, $h(e_1)=0$, so $dh+hd$ is its identity. Deleting the pair therefore preserves homology exactly. Repeating any sequence of presently valid unit cancellations remains correct by composition, even if sparsity gets worse. An acyclic discrete matching is a means of organizing such cancellations; no claim that an arbitrary proposed matching can be cancelled in arbitrary order is needed here.

\paragraph{A coefficient obstruction.}
Over $\mathbb Z$, the complex $\mathbb Z\xrightarrow{2}\mathbb Z$ has $H_0=\mathbb Z/2\mathbb Z$. Deleting its two generators because the entry is nonzero would instead leave zero homology. Over $\mathbb Q$ the entry is a unit and cancellation is valid. Thus integral computations need unit pivots, not merely nonzero incidences; stabilizer and orientation effects in a quotient complex must be accounted for before applying the rule.

\paragraph{A correspondence obstruction.}
Suppose a chain map $T$ is supplied together with the complex. If $i$ is the inclusion of the surviving complex in the changed bases and $p$ its projection, the transferred map is $pTi$. Both $i$ and $p$ are chain maps by the splitting, and the removed summand is contractible, so this formula induces the original map on homology. In contrast, coordinate deletion in the original bases can fail. Take
\[
 d:\mathbb Q\longrightarrow\mathbb Q^2,\quad d(1)=(1,1),\qquad
 T_1=2,\quad T_0=\begin{pmatrix}1&1\\0&2\end{pmatrix}.
\]
Here $T_0d=dT_1$, while on the surviving $H_0$ the correct projection and inclusion are $p(x,y)=y-x$ and $i(z)=(0,z)$. Thus $pT_0i=1$, whereas retaining only the original bottom-right entry would give $2$. Even a correct reduction of Betti numbers need not carry Hecke eigenvalues correctly without these maps.

\paragraph{The attempted efficiency inference stops.}
A single cancellation updates a potentially dense outer product $bu^{-1}a$. For example, with $u=1$, every entry of $a,b$ equal to one and $D=0$, it creates a full nonzero lower block $-ba$ from a sparse block matrix. This matrix is a legitimate two-term chain complex, since there are no adjacent nonzero differentials. Fewer generators alone therefore does not prove lower arithmetic or storage cost. A useful Voronoi matching must control fill-in, allowable coefficients, and transferred operators.
\section{Verification and adversarial audit}
The basis changes, adjacent differential conditions, contracting homotopy and transferred operator were checked explicitly. The unit requirement was tested with multiplication by two over both integers and rationals. The fill-in example is an actual chain complex but is not asserted to have a particular Voronoi realization.
\paragraph{Audit corrections.} Replaced nonzero-pivot cancellation by unit-pivot cancellation and included the Schur-complement and correspondence corrections. Distinguished abstract chain complexes from actual Voronoi incidence data.
\paragraph{External computation.} None; symbolic arguments and boundary checks only.
\section{Scope and remaining gaps}
No discrete matching on the original arithmetic complexes has been constructed, no performance bound is obtained, and torsion/stabilizer bookkeeping for those complexes remains untreated. The result is a fully justified local algebraic operation and concrete obstacles to two tempting shortcuts, not completion of the source's broad computational proposal.
\paragraph{Integration review.} Codex (GPT-6; exact serving revision not exposed) reviewed the written arguments and expanded layout and source-reference wrappers. Repository metadata and the production web renderer were checked using Python 3.13 and JavaScript with MathJax 3.2.2. This model-assisted review is not independent human proof verification; no PDF compilation or formal proof checking is asserted.

\section{Completion Estimate}
\noindent\textbf{COMPLETION: 0\%}
The calculations below the preserved source statement investigate a concrete route,
but establish no substantiated new progress toward the entire original target.
Known special cases and failed sufficient conditions are not credited as a solution.
\section{References}
The original source references and catalogue attribution are preserved in the
Sources section above. The following URLs were actually consulted for this attempt;
retention of an additional original reference does not claim fresh consultation of it.
\begin{itemize}
\item \url{https://ems.press/content/serial-article-files/46655} (consulted 2026-10-11). 
\end{itemize}
\section{Model and contributor disclosure}
OpenAI gpt-6-astra, requested reasoning configuration ultra, generated this research
attempt on 2026-10-11. The exact serving revision is not exposed. Codex GPT-6
performed orchestration, repository-format review, and written-proof review of this attempt; its exact serving revision
is also not exposed. The mathematical route, boundary checks and final adversarial
self-audit followed contribution prompts A then B in
\texttt{docs/contribution\_prompts.md}, after reading \texttt{CONTRIBUTING.md}.
Python 3.14.2 assembled the standalone documents; any mathematical computation is
specified in Verification. No theorem prover was used. No human mathematical
authorship, substantial human proof correction, or independent human proof checking
is asserted. Model-assisted review, source retrieval, and structural checks do not
establish correctness or novelty. Full private reasoning transcripts are unavailable
as reproducibility artifacts; the complete arguments, inputs and outcomes actually
used are stated here or in the companion batch audit.
\section{Final status}
Unresolved as a new research contribution. This is a concrete route investigation,
not a claim of new complete resolution or a current-status certification. A known
answer reproduced for context does not become a new contribution. The catalogue's
source-status label is preserved; this attempt status does not assert that a
source-labelled solved question is open.
\end{document}
